\subsection*{集合论}

可以说，在历史的每一个时期，数学家和哲学家都或多或少自觉地运用着集合论的论证。但在他们关于这一主题的观念史上，有必要将涉及基数概念（尤其是涉及无穷概念）的所有问题，与那些仅涉及属于和包含关系的问题明确区分开来。后者在直觉上是清晰的，似乎从未引起过争议。它们为三段论理论（如莱布尼茨和欧拉所展示的）、诸如“整体大于部分”这样的公理，以及几何学中涉及曲线与曲面相交的那部分内容，提供了最简便的基础。直到十九世纪末，对于谈论具有某一给定性质的所有对象的集合（或按某些作者的说法，“类”）(*)，同样没有人提出异议；而康托尔所给出的著名“定义”（“所谓集合，我们指的是将我们直觉或思维中的不同对象组合成一个整体”）({[}25{]}，第282页）在发表时几乎没有引起任何异议（†）。但当事关数或大小的概念与集合的概念纠缠在一起时，情况就大不相同了。长度无限可分的问题（当然是由早期毕达哥拉斯学派提出的）众所周知地导致了相当大的哲学困难；从埃利亚学派到博尔扎诺和康托尔，数学家和哲学家们都在由无限多个没有大小的点组成的有限大小这一悖论面前徒劳无功。对我们来说，追溯由这一问题引发的无休止且激烈的论战（即使只是概要）并不重要，这一问题构成了形而上学或神学谬误特别有利的土壤。我们只需注意到自古以来大多数数学家所持有的观点。它本质上在于拒绝讨论这个问题，因为它无法被无可辩驳地裁决，这种态度我们将在现代形式主义者中再次发现：正如后者设法消除所有“悖论性”集合的表象（见后文，第328页），古典数学家也小心翼翼地避免在他们的论证中引入“实无限”（即包含无限多个对象、被设想为同时存在（至少在思想中）的集合），而满足于“潜无限”，即任何给定大小（如果考虑的是“连续”大小，则是缩小它）增加的可能性（**）。如果这一观点包含了一定程度的虚伪（††），它仍然允许了古典数学的大部分发展（包括比例理论以及后来的无穷小演算）（*）；它似乎是一条极好的救命稻草，尤其是在无穷小引发的争吵之后，并且直到十九世纪中叶都几乎成为普遍接受的教条。

等势这一一般概念的初步萌芽出现在伽利略的一段评论中。{[}8{]}，第 VIII 卷，第 78-80 页）。他观察到该映射 \(n \to n^2\) 在自然整数与其平方之间建立了一一对应关系，并由此得出结论：公理“整体大于部分”不能应用于无限集合。但是，这一发现远未开启对无限集合的理性研究，似乎除了加深对“实无限”的不信任之外没有其他效果；这正是伽利略本人得出的教训，而柯西在1833年引用他时也只是为了赞同他的态度。

分析学的需求——尤其是对整个十九世纪都在深入研究的实变函数理论的更深入探讨——是后来成为现代集合论的东西的起源。当博尔扎诺在1817年证明实数集R中下有界集合的最大下界的存在性时，他像当时大多数同时代人一样，以“内涵”的方式论证，并且所考虑的并非实数的任意集合，而是它们的任意性质。但是，三十年后，当他撰写他的《无限悖论》时 {[}18{]} （出版于1851年，即他去世后三年），他毫不犹豫地主张“实无限”存在的可能性，并谈论任意集合。在这部著作中，他定义了集合等势的一般概念，并证明了R中任意两个紧区间都是等势的（前提是每个区间都包含多于一个点）。他还观察到有限集合与无限集合之间的特征性区别在于：无限集合E与E的某个真子集等势，但他没有给出这一断言的令人信服的证明。这部著作的整体基调更偏哲学而非数学；并且由于他没有足够精确地区分集合的势的概念与无穷大或无穷阶的大小概念，博尔扎诺在构造越来越大的势的无限集合的尝试中失败了，而他在这方面的讨论与对发散级数的考虑纠缠在一起，这些考虑完全没有意义。

集合论的创立，正如今天所理解的那样，归功于G.康托尔的天才。他也是从分析学家起步的，他在三角级数方面的工作，受黎曼的启发，在1872年自然地引导他首次尝试对这类理论中出现的“例外”集合进行分类（*），通过他在此背景下引入的逐次“导集”概念。毫无疑问，正是这些研究，以及他定义实数的方法，引导康托尔对等势问题产生兴趣；因为在1873年他注意到有理数集合（或代数数集合）是可数的。在他与戴德金大约此时开始的通信中 {[}25 bis{]}，他提出了整数集与实数集是否等势的问题，并在几周后成功证明了它们不等势。接着，从1874年起，维数问题一直萦绕在他心头，他花了三年时间徒劳地试图证明R与 \(R^{n}\) （n \textgreater{} 1）之间不可能存在一一对应，直到他最终构造出这样一个对应关系，连他自己都感到惊愕（†）。获得这些新颖而惊人的结果后，他全身心投入到集合论中。在1878年至1884年间发表于《数学年刊》的一系列六篇论文中，他考虑了等势问题、全序集理论、R与 \(R^{n}\) 的拓扑性质以及测度问题；令人钦佩的是，在他手中，那些看似与经典“连续统”概念密不可分的观念被逐步揭示出来。1880年，他萌生了“超限”迭代“导集”形成的想法；但这一想法直到两年后才随着良序集的引入而具体化——这是康托尔最具原创性的发现之一，凭借它，他得以对基数进行详细研究，并提出了“连续统问题”。 {[}25{]}.

如此大胆的构想，违背了两千年的传统，并导致了如此不可能且看似悖论的结果，指望它们毫无阻力地被接受是不现实的。事实上，在德国有影响力的数学家中，魏尔斯特拉斯是唯一对康托尔（他以前的学生）的工作表示某种支持的人；但康托尔遭到了施瓦茨，尤其是克罗内克（*）的不可调和反对。似乎正是反对他的思想所产生的持续紧张，以及他为证明连续统假设所做的徒劳努力，共同在康托尔身上引发了神经疾病的初期症状，这种疾病后来影响了他的数学产出（†）。事实上，直到大约1887年，他才重新对集合论产生兴趣，而他最后的出版物主要涉及全序集理论和序数运算，时间在1895-97年间。他还在1890年证明了不等式m \textless{} 2 \(^{m}\) 。然而，不仅连续统问题仍未得到解答，而且基数理论中还有一个更严重的缺陷，因为康托尔未能证明任意基数之间存在良序关系。这一缺陷部分由F.伯恩斯坦（1897年）的定理填补，该定理表明a ≤ b和b ≤ a蕴含a = b（**），而更重要的是由策梅洛的定理 {[}36 a{]} 填补，该定理证明了任何集合上都存在良序：这是康托尔早在1883年就已猜想过的定理（{[}25{]}，第169页）。

然而，戴德金从一开始就持续关注着康托尔的工作。但后者将注意力集中在无限集及其分类上，而戴德金则继续思考数的概念（这一思考已引导他通过“分割”定义了无理数）。在他1888年出版的小册子《Was sind und was sollen die Zahlen》中——其基本内容可追溯至1872-8年（{[}26{]}，第III卷，第335页）——他展示了自然整数的概念（正如我们所看到的，它如今是整个经典数学的基础）如何本身就能从集合论的基本概念中推导出来。他无疑是第一个明确发展一个集合到另一个集合的任意映射的基本性质的人（此前康托尔忽视了这一点，他只关心单射对应），并为集合E到自身的任意映射f引入了元素 \(a \in E\) 的“链”的概念，即所有集合 \(K \subset E\) 使得 \(a \in K\) 并且 \(f(K) \subset K (*)\) 的交集。接着，他将无限集E定义为存在一个E到E的一一映射 \(\varphi\) 使得 \(\varphi(E) \neq E (\dagger)\) 。此外，如果存在这样一个映射 \(\varphi\) 和一个元素 \(a \notin \varphi(E)\) 使得整个集合E是a的链，戴德金称E是“简单无限的”；而且，他注意到“皮亚诺公理”此时得到满足，并（在皮亚诺之前）展示了如何从这一出发点得到算术的所有基本定理。他的表述中唯一缺少的是无穷公理，戴德金（追随博尔扎诺）相信他可以通过将“思想世界”（Gedankenwelt）视为一个集合来证明它（**）。

在另一方面，戴德金因其算术工作（尤其是理想理论）而构想出比康托尔更为一般的有序集概念。后者完全局限于全序集（*），而戴德金则研究了一般情形，并特别对格进行了深入研究。{[}26{]}（第II卷，第236-271页）。这项工作在当时几乎没有引起任何关注；尽管他的结果后来被多位作者重新发现，并自1935年以来成为众多出版物的研究对象，但其历史重要性远不在于这一理论的可能应用（确实相当稀少），而在于它是精心公理化构造的首批范例之一。另一方面，康托尔关于可数集的早期结果很快产生了许多重要应用，甚至应用于最经典的分析问题（†）（更不用说他的工作中开创了一般拓扑学和测度论的部分；关于这些内容，读者可参阅《一般拓扑学》中的历史注释）。此外，十九世纪最后几年见证了超限归纳原理的首次应用，尤其是自策梅洛定理证明以来，该原理已成为现代数学所有部分不可或缺的工具。库拉托夫斯基在1922年给出了这一原理的一个版本，后来被佐恩重新发现 {[}45{]} 并且通常以他的名字命名，由于避免了良序集的使用，这一形式往往更便于应用（{[}40{]}，第89页）；如今普遍使用的正是这一形式的原理（**）。

因此，到十九世纪末，康托尔的基本观念已经占据了主导地位（††）。正如我们所看到的，数学的形式化在同一时期得以实现，公理化方法的使用也几乎被普遍接受。换言之，1875年至1895年间密集的研究工作为数学赢得了本书所呈现内容的本质部分。但同一时期也开启了一场罕见的“基础危机”，这场危机震撼了数学界三十余年，有时似乎不仅危及所有新近的成果，甚至危及数学中最经典的部分。

